\section{A square price is equivalent to planar flatness}
\label{sec:planar}

This section proves Theorem~\ref{thm:introplanar}, assuming
\cite[Proposition C.3, p.~28]{Anthropic2026}: every price is the
pullback of the metric of a complete $\cat$ space whose metric is
itself a price. The geometric rigidity below is classical
\cite[II.2.11]{BH1999}. We write out the polarization argument of
\cite[Lemma C.4(f), pp.~29, 65]{Anthropic2026}, based on the
quadrilateral inequality \cite{BN2008}.

\begin{lemma}[Flat rectangle]\label{lem:rectangle}
In a $\cat$ space suppose
\[
 d(x,y)=d(u,v)=a,\quad d(x,u)=d(y,v)=b,\quad
 d(x,v)^2=d(y,u)^2=a^2+b^2.
\]
If $p_s\in[x,u]$ and $q_t\in[y,v]$ occur at proportions
$s,t\in[0,1]$, then
$d(p_s,q_t)^2=a^2+b^2(s-t)^2$.
\end{lemma}

\begin{proof}
Define
\[
 B(p,p';r,r')=
 d(p,r')^2+d(p',r)^2-d(p,r)^2-d(p',r')^2.
\]
The quadrilateral inequality is
$B(p,p';r,r')\le2d(p,p')d(r,r')$.
Splitting both segments at $p_s,q_t$ expresses
$B(x,u;y,v)$ as a telescoping sum of four expressions.
Their upper bounds sum to $2b^2$, which is already
$B(x,u;y,v)$ by hypothesis. Hence each is an equality, including
\begin{equation}\label{eq:rectanglepolar}
 d(x,q_t)^2+d(y,p_s)^2-a^2-d(p_s,q_t)^2=2stb^2.
\end{equation}
For $t=s$, a second quadrilateral inequality bounds the first two
terms by $2s^2b^2+2a\,d(p_s,q_s)$. Thus
$(d(p_s,q_s)-a)^2\le0$ and $d(p_s,q_s)=a$.
Strong convexity of squared distance along a geodesic bounds each
of $d(x,q_s)^2$ and $d(y,p_s)^2$ by $a^2+s^2b^2$.
Their sum attains that bound, so both equalities hold.
Renaming $s$ as $t$ in the first and substituting into
\eqref{eq:rectanglepolar} gives the claim.
\end{proof}

\begin{proof}[Proof of Theorem~\ref{thm:introplanar}]
The third assertion implies the second, which implies the first
by restriction. Embed a square price in a price space and name
its vertices $x,y,v,u$ in cyclic order.
Let $l_t,r_t$ be proportion-$t$ points on its left and right
vertical edges. The lemma gives
\[
 d(l_t,r_{t'})^2=1+(t-t')^2,\qquad
 d(l_t,l_{t'})=d(r_t,r_{t'})=|t-t'|.
\]
Let $F(s,t)$ be the proportion-$s$ point on $[l_t,r_t]$.
For $t\ne t'$, the four endpoints
$l_t,r_t,r_{t'},l_{t'}$ form a rectangle with sides $1$ and
$|t-t'|$. Apply the lemma along its horizontal edges:
\[
 d(F(s,t),F(s',t'))^2=(s-s')^2+(t-t')^2.
\]
When $t=t'$, the same formula is the constant-speed geodesic
formula. Therefore the whole unit square embeds isometrically.
Every finite planar set can be translated and scaled into it.
Pullback and rescaling preserve prices, proving the third assertion.
The same construction fills any rectangle with positive area,
and every finite planar set fits in a homothetic copy of it.
\end{proof}

\begin{corollary}\label{cor:planarobstruction}
Assuming the price-space embedding theorem, the existence of any
finite planar Euclidean configuration that is not a price is
equivalent to the square not being a price.
If \cite[Theorem D.2, pp.~30--31]{Anthropic2026} also holds,
either failure implies the contracting hypothesis.
\end{corollary}

\begin{proof}
The equivalence is the contrapositive of
Theorem~\ref{thm:introplanar}. A failed planar instance refutes
the full flat hypothesis, so Theorem D.2 gives the contracting one.
\end{proof}

\subsection{Cross-polytopes detect flatness in every fixed dimension}

The four vertices of a square are the two-dimensional instance of a
more general rigid configuration. Write
\[
 X_D=\{e_1,-e_1,\ldots,e_D,-e_D\}\subset\mathbb R^D,
 \qquad C_D=\operatorname{conv}(X_D)
           =\{x:\textstyle\sum_i|x_i|\le1\}.
\]
The next statement is classical metric geometry in the sense that its
inputs are the flat-square theorem, the parallel-set splitting theorem,
the tangent-cone theorem and the Kirszbraun--Lang--Schroeder extension
theorem. We give the argument because a collection of pairwise flat
coordinate squares does not, without an additional argument, justify
the existence of a higher-dimensional flat.

\begin{lemma}[Cross-polytope rigidity]\label{lem:cross-polytope-flat}
Let $D\ge1$ and let $f:X_D\longrightarrow X$ be an isometric embedding
into a complete $\cat$ space. There is an isometric embedding
$F:C_D\longrightarrow X$ extending $f$. Its image is precisely the
closed convex hull of $f(X_D)$. In particular that hull contains an
isometric Euclidean $D$-ball of radius $1/\sqrt D$.
\end{lemma}

\begin{proof}
For $D=1$ use the geodesic between the two vertices. Suppose $D\ge2$.
Every two antipodal pairs have the metric of a Euclidean square, so
Lemma~\ref{lem:rectangle} and its filling argument give an isometric
filled square. The two diagonals in each such square have the same
midpoint. Consequently all $D$ antipodal segments have one common
midpoint, denoted by $o$.

Let $T$ be the completed tangent cone of $X$ at $o$. It is a complete
$\cat$ space by \cite[II.3.19, p.~191]{BH1999}. The logarithm map
$\log_o:X\longrightarrow T$, sending a point to its distance and
initial direction from $o$, is $1$-Lipschitz: comparison angles at $o$
are at least the Alexandrov angle, which gives
\[
 d_T(\log_o x,\log_o y)^2
 \le d_X(x,y)^2.
\]
On each filled square through $o$ it preserves Euclidean distances.
The cone over the directions of that square therefore contains a full
Euclidean plane. It follows that $T$ contains $D$ geodesic lines
$L_1,\ldots,L_D$ through its vertex, with each pair forming the
coordinate axes of an isometrically embedded Euclidean plane.

We now justify that these lines lie in one Euclidean $D$-flat. More
generally, this is true in any complete $\cat$ space with the stated
configuration of complete lines. Induct on $D$. The parallel set of
$L_1$ is convex and splits as $Y\times\mathbb R$, where the second
coordinate is projection onto $L_1$;
see \cite[II.2.14, p.~183]{BH1999}. If necessary pass to its closure,
which is the product of the completion of $Y$ with $\mathbb R$; this
closure remains a convex subspace of the given complete space. The
plane formed by $L_1$ and $L_j$ shows that $L_j$ is contained in this
parallel set and has second coordinate zero. For $j,k>1$, the entire
plane formed by $L_j$ and $L_k$ is the convex hull of those two entire
lines: a point with coordinates $(a,b)$ is the midpoint of $(2a,0)$
and $(0,2b)$. Hence this plane is in the parallel set. Its
$\mathbb R$ coordinate is zero because that coordinate is affine on
geodesics. The completed factor $Y$ thus contains $D-1$ lines with
the same pairwise-plane property. Induction gives a Euclidean
$(D-1)$-flat in that factor, and its product with $\mathbb R$ gives
the required $D$-flat. No local compactness or geodesic completeness
of $X$ has been assumed.

Identify this flat with $\mathbb R^D$ so that
$\log_o f(\pm e_i)=\pm e_i$. In particular $C_D$ is a closed convex
subset of $T$. The Kirszbraun--Lang--Schroeder theorem
\cite{LangSchroeder1997}, in the form
\cite[Theorem~10.14, p.~160]{AKP2023}, extends the vertex map $f$ to a
$1$-Lipschitz map $F:\mathbb R^D\longrightarrow X$. Here the source
is complete with curvature bounded below by zero and the target is
complete with curvature bounded above by zero, so the extension
theorem has no radius restriction. Restrict $F$ to $C_D$ and let
$P:T\longrightarrow C_D$ be nearest-point projection. Projection
onto a complete convex subset of a $\cat$ space is $1$-Lipschitz
\cite[II.2.4]{BH1999}. Thus
\[
 G=P\circ\log_o\circ F:C_D\longrightarrow C_D
\]
is $1$-Lipschitz and fixes every vertex.

The fixed set of a $1$-Lipschitz self-map of a uniquely geodesic space
is convex. Indeed, if $a,b$ are fixed and $p\in[a,b]$, the inequalities
$d(a,Gp)\le d(a,p)$ and $d(Gp,b)\le d(p,b)$ have total upper bound
$d(a,b)$, so both are equalities and uniqueness of the geodesic gives
$Gp=p$. Therefore $G$ fixes the convex hull of its vertices, namely
all of $C_D$. For $p,q\in C_D$ we consequently have
\[
 \|p-q\|=d_T(Gp,Gq)
 \le d_T(\log_o Fp,\log_o Fq)
 \le d_X(Fp,Fq)\le\|p-q\|.
\]
Every inequality is an equality, proving that $F|_{C_D}$ is isometric.
Its image is convex by uniqueness of geodesics, is closed by
compactness, and is generated by the images of the vertices under
finite convex combinations. It is therefore their closed convex
hull. Finally $\sum_i|x_i|\le\sqrt D\|x\|$ shows that $C_D$ contains
the claimed ball.
\end{proof}

A recent, stronger variance-rigidity statement also implies this
lemma: \cite[Proposition~1.4 and Section~13.1]{Lytchak2026geometry}
states that equality in the $\cat$ variance inequality forces the
closed convex hull of the measure's support to be a convex subset of
a Hilbert space. For the uniform measure on $f(X_D)$, the midpoint
inequality for each antipodal pair gives
\[
 \frac1{2D}\sum_{v\in X_D}d(z,f(v))^2\ge1+d(z,o)^2,
 \qquad
 \frac1{(2D)^2}\sum_{u,v\in X_D}d(f(u),f(v))^2=2.
\]
Thus $o$ is its barycentre with variance $1$, and equality holds in
that criterion. The preceding proof uses the classical extension and
splitting theorems directly and does not require the newer preprint.

\begin{theorem}[A finite test for flatness in dimension $D$]
\label{thm:cross-polytope-price-flat}
Fix a field and $D\ge1$. Assume
\cite[Proposition~C.3, p.~28]{Anthropic2026}, including that every
finite subset of its complete $\cat$ price space is a price. The
following conditions are equivalent:
\begin{enumerate}
\item the Euclidean metric of $X_D$ is a price;
\item the Euclidean metric of $X_D\cup\{0\}$ is a price;
\item every finite subset of $\mathbb R^D$ has its Euclidean metric
as a price.
\end{enumerate}
\end{theorem}

\begin{proof}
The third condition implies the second and the second implies the
first by restriction. Under the first condition, embed $X_D$ into a
complete $\cat$ price space using Proposition~C.3. The embedding is
injective because the vertex distances are positive. By
Lemma~\ref{lem:cross-polytope-flat}, this space contains a Euclidean
$D$-ball of positive radius. Any finite subset of $\mathbb R^D$ can
be translated and scaled into that ball. Its image has a price
metric, and pullback followed by inverse rescaling preserves prices.
This proves the third condition.
\end{proof}

In dimension two, $X_2$ is a rotated square of side $\sqrt2$;
Theorem~\ref{thm:cross-polytope-price-flat} therefore recovers
Theorem~\ref{thm:introplanar}. In higher dimensions it supplies the
geometric step needed to compare a shortening of the $D$ coordinate
diameters with failure of flatness in that same dimension.
